\documentclass[aspectratio=169, 14pt]{beamer}
\usepackage{beamer_preamble}

\begin{document}

% ==== Title ===============================================================
\titleslide{Lesson 4-3 \textperiodcentered\ Period 2}{Multiply \& Divide Rational Expressions + DOK-3 Closure}

% ==== Do Now ==============================================================
\begin{frame}{Do Now --- Simplify and state the domain}{Algebra 2 \textperiodcentered\ Lesson 4-3 \textperiodcentered\ Period 2}
\phasebadges{1}{5}{bridge from P1}
\vspace{0.2cm}
\begin{projcard}{Do Now \textperiodcentered\ Simplify this rational expression}{slidegold}
What is the simplified form of the rational expression? What is the domain?

\vspace{0.25cm}
\[
\frac{x^3+9x^2-10x}{x^3-9x^2-10x}
\]

\vspace{0.1cm}
\small Factor every numerator and denominator completely. State all excluded values.
\end{projcard}
\end{frame}

% ==== Launch ==============================================================
\begin{frame}{Launch --- Example 3 (Multiply \& Divide Rules)}{Algebra 2 \textperiodcentered\ Lesson 4-3 \textperiodcentered\ Period 2}
\phasebadges{2}{12}{teacher models}
\vspace{0.2cm}
\begin{projcard}{Multiply / Divide Rules (keep on your page)}{slidegreen}
\begin{enumerate}
  \setlength{\itemsep}{1pt}
  \item Factor every numerator AND denominator completely.
  \item For division, rewrite as multiplication by the \textbf{RECIPROCAL} of the divisor.
  \item Cancel only common \textbf{FACTORS} across any numerator with any denominator.
  \item State the domain: exclude every zero of every denominator that appeared in \textbf{ANY} step, including the reciprocal's new denominator.
\end{enumerate}

\vspace{0.1cm}
\small
Division rule: \(\displaystyle\frac{a}{b}\div\frac{c}{d}=\frac{a}{b}\times\frac{d}{c}\)
\hspace{1em}
\textcolor{slideaccent}{\textbf{Flip FIRST --- then cancel.}}
\end{projcard}
\end{frame}

% ==== Explore =============================================================
\begin{frame}{Explore --- TPS + DOK-3 Performance Task}{Algebra 2 \textperiodcentered\ Lesson 4-3 \textperiodcentered\ Period 2}
\phasebadges{2-3}{33}{student work}
\small\textcolor{slidegray}{7 items in order. Plan \(\sim\)12 min for Practice \#13 (Closure Argument).}

\vspace{0.10cm}
\begin{projcard}{Try It 3 \textperiodcentered\ Multiply rational expressions}{slideblue}
Factor, cancel, state domain. (Two parts.)
\end{projcard}
\vspace{-0.12cm}
\begin{projcard}{Practice \#27 \textperiodcentered\ Multiply with factoring}{slideblue}
\(\displaystyle\frac{(x-y)^2}{x+y}\cdot\frac{3x+3y}{x^2-y^2}\) --- find product and domain.
\end{projcard}
\vspace{-0.12cm}
\begin{projcard}{Try It 4 \textperiodcentered\ Divide rational expressions}{slideblue}
Flip the divisor first. State domain including reciprocal's denominator.
\end{projcard}
\vspace{-0.12cm}
\begin{projcard}{Practice \#29 \textperiodcentered\ Divide: rewrite as multiplication}{slideblue}
\(\displaystyle\frac{2x^2-10x}{(x-5)(x^2-1)}\cdot(3x^2+4x+1)\) --- domain?
\end{projcard}
\vspace{-0.12cm}
\begin{projcard}{Try It 5 + Practice \#32 \textperiodcentered\ Mixed}{slideblue}
Mixed multiply/divide, multi-step. State full domain.
\end{projcard}
\vspace{-0.12cm}
\begin{projcard}{Practice \#13 \textperiodcentered\ PERFORMANCE TASK \(\bigl(\bigstar\bigr)\)}{slideaccent}
\textbf{DOK-3 Closure Argument} --- Why are rational expressions closed under multiplication and division? Write a paragraph with ONE product example AND one quotient example.
\end{projcard}
\end{frame}

% ==== Share / Summary =====================================================
\begin{frame}{Share / Summary}{Algebra 2 \textperiodcentered\ Lesson 4-3 \textperiodcentered\ Period 2}
\phasebadges{2}{5}{academic conversation + P3 bridge}
\vspace{0.2cm}
\begin{projcard}{Sentence frame \textperiodcentered\ Closure Argument}{slideblue}
``I rewrite the division as multiplication by the reciprocal, giving \underline{\hspace{1cm}}. After factoring, I cancel the factor \underline{\hspace{1cm}}, leaving \underline{\hspace{1cm}}. The domain excludes \(x =\) \underline{\hspace{1cm}} because those values make at least one denominator zero at some step.''
\end{projcard}

\vspace{0.15cm}
\begin{projcard}{Bridge to Period 3}{slideaccent}
Next period we apply multiply/divide skills to complex rational expressions and assessment-critical items. Bring your domain-restriction list habit.
\end{projcard}
\end{frame}

% ==== Exit Ticket =========================================================
\begin{frame}{Exit Ticket --- Summary Recap}{Algebra 2 \textperiodcentered\ Lesson 4-3 \textperiodcentered\ Period 2}
\phasebadges{1-2}{3}{not graded}
\vspace{0.2cm}
\begin{projcard}{Complete on your exit ticket}{slidegold}
\begin{enumerate}
  \setlength{\itemsep}{0.2cm}
  \item To divide rational expressions I first \underline{\hspace{5cm}}.
  \item I cancel common factors, not \underline{\hspace{5cm}}.
  \item My domain restriction list must include zeros from \underline{\hspace{5cm}}.
  \item One biggest thing I learned today: \underline{\hspace{4.5cm}}.
\end{enumerate}
\end{projcard}
\end{frame}

\end{document}
